\chapter{Cardioversion}

\section*{What Is This Test or Procedure?}
Cardioversion restores a normal heart rhythm in people with certain types of abnormal heartbeats (arrhythmias), usually by delivering a synchronized electrical shock through electrode pads on the chest.

\section*{Alternative Names}
Electrical cardioversion, Direct current cardioversion, DCCV

\section*{Principle}
A synchronized electrical shock delivered through pads or paddles depolarizes the heart muscle simultaneously, abruptly terminating reentrant electrical circuits. The pause lets the sinoatrial node resume control and a coordinated, normal rhythm be restored.
\section*{History}
Bernard Lown developed the direct current cardioverter-defibrillator in 1961, introducing synchronized cardioversion for atrial fibrillation and flutter, which was much safer than chemical methods.

\section*{Why This Test or Procedure Is Ordered}
Ordered to treat atrial fibrillation or atrial flutter when medications fail, or when the patient is hemodynamically unstable (e.g., severe hypotension, pulmonary edema).

\section*{Is This a Primary or Secondary Test or Procedure}
Primary treatment for unstable tachyarrhythmias; elective treatment for persistent symptomatic atrial fibrillation.

\section*{How This Test or Procedure Is Performed}
Performed under brief general anesthesia or heavy sedation. Electrode pads are placed on the chest (anterior-posterior or anterior-lateral). The defibrillator delivers a shock synchronized to the R wave of the QRS complex to avoid triggering ventricular fibrillation.

\section*{What Preparation Is Needed}
Fast for 6-8 hours. Ensure adequate anticoagulation for at least 3 weeks prior, or perform a transesophageal echocardiogram (TEE) to rule out left atrial thrombus.

\section*{Contraindications and Precautions}
Known left atrial thrombus, digitalis toxicity-induced arrhythmia, or uncorrected hypokalemia.

\section*{Risks}
Thromboembolism (stroke), transient arrhythmias, skin irritation/burns from pads, or respiratory depression from sedation.

\section*{What Happens After the Test or Procedure}
Monitor vital signs and airway until fully awake. Verify rhythm with a 12-lead ECG. Continue oral anticoagulation as directed.

\section*{Understanding Your Results}
Conversion of the arrhythmia to sinus rhythm is immediately visible on the monitor.

\section*{Normal Results}
Restoration and maintenance of normal sinus rhythm.

\section*{Follow-up Tests and Procedures}
ECG check-up; maintenance of antiarrhythmic and anticoagulant therapy as indicated.

\section*{References}
\begin{enumerate}
  \item MedlinePlus. Cardioversion. U.S. National Library of Medicine; 2025.
  \item Current Medical Diagnosis and Treatment. McGraw-Hill; 2025.
\end{enumerate}

\clearpage
